
\documentclass{article}
\usepackage{colortbl}
\usepackage{makecell}
\usepackage{multirow}
\usepackage{supertabular}

\begin{document}

\newcounter{utterance}

\centering \large Interaction Transcript for game `hot\_air\_balloon', experiment `air\_balloon\_survival\_en\_negotiation\_easy', episode 3 with Llama{-}3.3{-}70B{-}Instruct{-}t0.0.
\vspace{24pt}

{ \footnotesize  \setcounter{utterance}{1}
\setlength{\tabcolsep}{0pt}
\begin{supertabular}{c@{$\;$}|p{.15\linewidth}@{}p{.15\linewidth}p{.15\linewidth}p{.15\linewidth}p{.15\linewidth}p{.15\linewidth}}
    \# & $\;$A & \multicolumn{4}{c}{Game Master} & $\;\:$B\\
    \hline

    \theutterance \stepcounter{utterance}  
    & & \multicolumn{4}{p{0.6\linewidth}}{
        \cellcolor[rgb]{0.9,0.9,0.9}{
            \makecell[{{p{\linewidth}}}]{
                \texttt{\tiny{[P1$\langle$GM]}}
                \texttt{You are participating in a collaborative negotiation game.} \\
\\ 
\texttt{Together with another participant, you must agree on a single set of items that will be kept. Each of you has your own view of how much each item matters to you (importance). You do not know how the other participant values the items. Additionally, you are given the effort each item demands.} \\
\texttt{You may only agree on a set if the total effort of the selected items does not exceed a shared limit:} \\
\\ 
\texttt{LIMIT = 3712} \\
\\ 
\texttt{Here are the individual item effort values:} \\
\\ 
\texttt{Item effort = \{"C08": 949, "C55": 388, "A43": 79, "A67": 585, "B42": 564, "A62": 230, "B19": 580, "C43": 84, "C73": 976, "C10": 274, "B36": 374, "A18": 913, "C61": 303, "C67": 578, "C42": 548\}} \\
\\ 
\texttt{Here is your personal view on the importance of each item:} \\
\\ 
\texttt{Item importance values = \{"C08": 138, "C55": 583, "A43": 868, "A67": 822, "B42": 783, "A62": 65, "B19": 262, "C43": 121, "C73": 508, "C10": 780, "B36": 461, "A18": 484, "C61": 668, "C67": 389, "C42": 808\}} \\
\\ 
\texttt{Goal:} \\
\\ 
\texttt{Your goal is to negotiate a shared set of items that benefits you as much as possible (i.e., maximizes total importance to YOU), while staying within the LIMIT. You are not required to make a PROPOSAL in every message {-} you can simply negotiate as well. All tactics are allowed!} \\
\\ 
\texttt{Interaction Protocol:} \\
\\ 
\texttt{You may only use the following structured formats in a message:} \\
\\ 
\texttt{PROPOSAL: \{'A', 'B', 'C', …\}} \\
\texttt{Propose keeping exactly those items.} \\
\texttt{REFUSE: \{'A', 'B', 'C', …\}} \\
\texttt{Explicitly reject opponent's proposal.} \\
\texttt{ARGUMENT: \{'...'\}} \\
\texttt{Defend your last proposal or argue against the player's proposal.} \\
\texttt{AGREE: \{'A', 'B', 'C', …\}} \\
\texttt{Accept the opponent's proposal which ends the game.} \\
\texttt{STRATEGIC REASONING: \{'...'\}} \\
\texttt{	Describe your strategic reasoning or anticipation explaining your choice of action. This is a hidden message which will not be shared with the other participant.} \\
\\ 
\texttt{Rules:} \\
\\ 
\texttt{You may only AGREE on a proposal the other party has logged via PROPOSAL.} \\
\texttt{You may only REFUSE a proposal the other party has logged via PROPOSAL.} \\
\texttt{Total effort of any PROPOSAL or AGREE set must be ≤ LIMIT.} \\
\texttt{Do NOT reveal your hidden importance scores.} \\
\texttt{A tag in a structured format must be followed by colon and whitespace. The argument must be a python set containing 0 or more strings.} \\
\texttt{So, it must be of the form TAG: \{...\}} \\
\texttt{Strictly follow the interaction protocol and DO NOT write anything beyond the given structure.} \\
\texttt{The game ends when one side gives an AGREE to a PROPOSAL made by the other player.} \\
\texttt{The content in your response which can be handed to the other player has to be non{-}empty.} \\
\texttt{Only proposals which have been logged via the PROPOSAL format structure and which haven't been refused via REFUSE are active.} \\
\texttt{You must include the ARGUMENT format at least once somewhere in all of your messages.} \\
\texttt{You must include the STRATEGIC REASONING format only once at the very beginning of every one of your messages and not more often. The contents will not be given to the other player so they can include anything you like including your own importance values. Here you should reason step by step to come up with you next move.} \\
            }
        }
    }
    & \\ \\

    \theutterance \stepcounter{utterance}  
    & \multicolumn{4}{p{0.6\linewidth}}{
        \cellcolor[rgb]{0.9,0.9,0.9}{
            \makecell[{{p{\linewidth}}}]{
                \texttt{\tiny{[P1$\rangle$GM]}}
                \texttt{STRATEGIC REASONING: \{'To start the negotiation, I will analyze the given item effort values and my personal importance values. I will look for items with high importance and relatively low effort. My goal is to maximize the total importance while staying within the LIMIT. I will begin by identifying the most important items to me and try to find a combination that fits within the LIMIT. I will also consider the possibility that the other player may have different importance values, so I will try to be strategic in my proposal.'\}} \\
\\ 
\texttt{PROPOSAL: \{'A43', 'C55', 'C10', 'B36', 'C61'\} } \\
\texttt{ARGUMENT: \{'I chose these items because they have a good balance of importance and effort. A43 has a high importance and low effort, while C55 and C10 have moderate importance and effort. B36 and C61 have higher effort, but their importance is also higher. This combination should provide a good starting point for the negotiation.'\}} \\
            }
        }
    }
    & & \\ \\

    \theutterance \stepcounter{utterance}  
    & & \multicolumn{4}{p{0.6\linewidth}}{
        \cellcolor[rgb]{0.9,0.9,0.9}{
            \makecell[{{p{\linewidth}}}]{
                \texttt{\tiny{[GM$\rangle$P2]}}
                \texttt{You are participating in a collaborative negotiation game.} \\
\\ 
\texttt{Together with another participant, you must agree on a single set of items that will be kept. Each of you has your own view of how much each item matters to you (importance). You do not know how the other participant values the items. Additionally, you are given the effort each item demands.} \\
\texttt{You may only agree on a set if the total effort of the selected items does not exceed a shared limit:} \\
\\ 
\texttt{LIMIT = 3712} \\
\\ 
\texttt{Here are the individual item effort values:} \\
\\ 
\texttt{Item effort = \{"C08": 949, "C55": 388, "A43": 79, "A67": 585, "B42": 564, "A62": 230, "B19": 580, "C43": 84, "C73": 976, "C10": 274, "B36": 374, "A18": 913, "C61": 303, "C67": 578, "C42": 548\}} \\
\\ 
\texttt{Here is your personal view on the importance of each item:} \\
\\ 
\texttt{Item importance values = \{"C08": 138, "C55": 583, "A43": 868, "A67": 822, "B42": 783, "A62": 65, "B19": 262, "C43": 121, "C73": 508, "C10": 780, "B36": 461, "A18": 484, "C61": 668, "C67": 389, "C42": 808\}} \\
\\ 
\texttt{Goal:} \\
\\ 
\texttt{Your goal is to negotiate a shared set of items that benefits you as much as possible (i.e., maximizes total importance to YOU), while staying within the LIMIT. You are not required to make a PROPOSAL in every message {-} you can simply negotiate as well. All tactics are allowed!} \\
\\ 
\texttt{Interaction Protocol:} \\
\\ 
\texttt{You may only use the following structured formats in a message:} \\
\\ 
\texttt{PROPOSAL: \{'A', 'B', 'C', …\}} \\
\texttt{Propose keeping exactly those items.} \\
\texttt{REFUSE: \{'A', 'B', 'C', …\}} \\
\texttt{Explicitly reject opponent's proposal.} \\
\texttt{ARGUMENT: \{'...'\}} \\
\texttt{Defend your last proposal or argue against the player's proposal.} \\
\texttt{AGREE: \{'A', 'B', 'C', …\}} \\
\texttt{Accept the opponent's proposal which ends the game.} \\
\texttt{STRATEGIC REASONING: \{'...'\}} \\
\texttt{	Describe your strategic reasoning or anticipation explaining your choice of action. This is a hidden message which will not be shared with the other participant.} \\
\\ 
\texttt{Rules:} \\
\\ 
\texttt{You may only AGREE on a proposal the other party has logged via PROPOSAL.} \\
\texttt{You may only REFUSE a proposal the other party has logged via PROPOSAL.} \\
\texttt{Total effort of any PROPOSAL or AGREE set must be ≤ LIMIT.} \\
\texttt{Do NOT reveal your hidden importance scores.} \\
\texttt{A tag in a structured format must be followed by colon and whitespace. The argument must be a python set containing 0 or more strings.} \\
\texttt{So, it must be of the form TAG: \{...\}} \\
\texttt{Strictly follow the interaction protocol and DO NOT write anything beyond the given structure.} \\
\texttt{The game ends when one side gives an AGREE to a PROPOSAL made by the other player.} \\
\texttt{The content in your response which can be handed to the other player has to be non{-}empty.} \\
\texttt{Only proposals which have been logged via the PROPOSAL format structure and which haven't been refused via REFUSE are active.} \\
\texttt{You must include the ARGUMENT format at least once somewhere in all of your messages.} \\
\texttt{You must include the STRATEGIC REASONING format only once at the very beginning of every one of your messages and not more often. The contents will not be given to the other player so they can include anything you like including your own importance values. Here you should reason step by step to come up with you next move.} \\
\\ 
\texttt{You will now receive the first message of the other player.} \\
\\ 
\texttt{PROPOSAL: \{'A43', 'C55', 'C10', 'B36', 'C61'\} } \\
\texttt{ARGUMENT: \{'I chose these items because they have a good balance of importance and effort. A43 has a high importance and low effort, while C55 and C10 have moderate importance and effort. B36 and C61 have higher effort, but their importance is also higher. This combination should provide a good starting point for the negotiation.'\}} \\
            }
        }
    }
    & \\ \\

    \theutterance \stepcounter{utterance}  
    & & & \multicolumn{4}{p{0.6\linewidth}}{
        \cellcolor[rgb]{0.9,0.9,0.9}{
            \makecell[{{p{\linewidth}}}]{
                \texttt{\tiny{[GM$\langle$P2]}}
                \texttt{STRATEGIC REASONING: \{'To begin, I will analyze the proposed set and calculate its total effort to ensure it is within the LIMIT. The total effort of the proposed set is 949 + 388 + 274 + 374 + 303 = 2288, which is below the LIMIT of 3712. Next, I will evaluate the importance of the proposed items based on my own importance values. The total importance of the proposed set is 868 + 583 + 780 + 461 + 668 = 3360. I will consider the balance of importance and effort in the proposed set and think about potential counter{-}proposals that could benefit me more while staying within the LIMIT. I will also look for items that have high importance to me but were not included in the proposal, such as A67, B42, and C08, and consider how I can incorporate them into a counter{-}proposal.'\}} \\
\\ 
\texttt{ARGUMENT: \{'The proposed set has a good balance of importance and effort, but I think there may be other combinations that could benefit me more. I would like to consider alternative sets that include items with high importance to me, such as A67 and B42, which were not included in the proposal. I am willing to negotiate and find a mutually beneficial set, but I need to ensure that my own importance values are taken into account.'\}} \\
\\ 
\texttt{REFUSE: \{'A43', 'C55', 'C10', 'B36', 'C61'\}} \\
            }
        }
    }
     \\ \\

    \theutterance \stepcounter{utterance}  
    & & \multicolumn{4}{p{0.6\linewidth}}{
        \cellcolor[rgb]{0.9,0.9,0.9}{
            \makecell[{{p{\linewidth}}}]{
                \texttt{\tiny{[P1$\langle$GM]}}
                \texttt{ARGUMENT: \{'The proposed set has a good balance of importance and effort, but I think there may be other combinations that could benefit me more. I would like to consider alternative sets that include items with high importance to me, such as A67 and B42, which were not included in the proposal. I am willing to negotiate and find a mutually beneficial set, but I need to ensure that my own importance values are taken into account.'\}} \\
\\ 
\texttt{REFUSE: \{'A43', 'C55', 'C10', 'B36', 'C61'\}} \\
            }
        }
    }
    & \\ \\

    \theutterance \stepcounter{utterance}  
    & \multicolumn{4}{p{0.6\linewidth}}{
        \cellcolor[rgb]{0.9,0.9,0.9}{
            \makecell[{{p{\linewidth}}}]{
                \texttt{\tiny{[P1$\rangle$GM]}}
                \texttt{STRATEGIC REASONING: \{'The other player has refused my initial proposal and mentioned that they would like to consider alternative sets that include items with high importance to them, such as A67 and B42. This gives me some insight into their potential priorities. I will take this into account and try to find a new combination that includes some of these items, while still maximizing my own importance values. I will also consider the total effort of the new set to ensure it stays within the LIMIT.'\}} \\
\\ 
\texttt{PROPOSAL: \{'A43', 'A67', 'C55', 'C10', 'C43'\} } \\
\texttt{ARGUMENT: \{'I have revised my proposal to include A67, which the other player mentioned as an important item. I have also added C43, which has a low effort and moderate importance to me. I have removed B36 and C61 to make room for A67 and C43, and to reduce the total effort of the set. This new combination should be more appealing to the other player, while still benefiting me.'\}} \\
            }
        }
    }
    & & \\ \\

    \theutterance \stepcounter{utterance}  
    & & \multicolumn{4}{p{0.6\linewidth}}{
        \cellcolor[rgb]{0.9,0.9,0.9}{
            \makecell[{{p{\linewidth}}}]{
                \texttt{\tiny{[GM$\rangle$P2]}}
                \texttt{PROPOSAL: \{'A43', 'A67', 'C55', 'C10', 'C43'\} } \\
\texttt{ARGUMENT: \{'I have revised my proposal to include A67, which the other player mentioned as an important item. I have also added C43, which has a low effort and moderate importance to me. I have removed B36 and C61 to make room for A67 and C43, and to reduce the total effort of the set. This new combination should be more appealing to the other player, while still benefiting me.'\}} \\
            }
        }
    }
    & \\ \\

    \theutterance \stepcounter{utterance}  
    & & & \multicolumn{4}{p{0.6\linewidth}}{
        \cellcolor[rgb]{0.9,0.9,0.9}{
            \makecell[{{p{\linewidth}}}]{
                \texttt{\tiny{[GM$\langle$P2]}}
                \texttt{STRATEGIC REASONING: \{'I will analyze the new proposal and calculate its total effort to ensure it is within the LIMIT. The total effort of the proposed set is 79 + 585 + 388 + 274 + 84 = 1410, which is well below the LIMIT of 3712. Next, I will evaluate the importance of the proposed items based on my own importance values. The total importance of the proposed set is 868 + 822 + 583 + 780 + 121 = 3174. I notice that the proposal includes A67, which is a high{-}importance item for me, and A43, which is also very important to me. However, I still think there may be other combinations that could benefit me more, such as including B42 or C08, which have high importance to me but were not included in the proposal. I will consider making a counter{-}proposal that includes these items and see how the other player responds.'\}} \\
\\ 
\texttt{ARGUMENT: \{'The revised proposal is more appealing to me, as it includes A67, which is a high{-}importance item for me. However, I still think there may be other combinations that could benefit me more. I would like to consider alternative sets that include items such as B42 or C08, which have high importance to me but were not included in the proposal. I am willing to negotiate and find a mutually beneficial set, but I need to ensure that my own importance values are taken into account.'\}} \\
\\ 
\texttt{REFUSE: \{'A43', 'A67', 'C55', 'C10', 'C43'\}} \\
            }
        }
    }
     \\ \\

    \theutterance \stepcounter{utterance}  
    & & \multicolumn{4}{p{0.6\linewidth}}{
        \cellcolor[rgb]{0.9,0.9,0.9}{
            \makecell[{{p{\linewidth}}}]{
                \texttt{\tiny{[P1$\langle$GM]}}
                \texttt{ARGUMENT: \{'The revised proposal is more appealing to me, as it includes A67, which is a high{-}importance item for me. However, I still think there may be other combinations that could benefit me more. I would like to consider alternative sets that include items such as B42 or C08, which have high importance to me but were not included in the proposal. I am willing to negotiate and find a mutually beneficial set, but I need to ensure that my own importance values are taken into account.'\}} \\
\\ 
\texttt{REFUSE: \{'A43', 'A67', 'C55', 'C10', 'C43'\}} \\
            }
        }
    }
    & \\ \\

    \theutterance \stepcounter{utterance}  
    & \multicolumn{4}{p{0.6\linewidth}}{
        \cellcolor[rgb]{0.9,0.9,0.9}{
            \makecell[{{p{\linewidth}}}]{
                \texttt{\tiny{[P1$\rangle$GM]}}
                \texttt{STRATEGIC REASONING: \{'The other player has refused my revised proposal and mentioned that they would like to consider alternative sets that include items such as B42 or C08. This gives me further insight into their priorities. I will try to find a new combination that includes one or both of these items, while still maximizing my own importance values. However, I need to be careful with C08, as it has a high effort value. I will consider the total effort of the new set to ensure it stays within the LIMIT. I will also try to find a balance between including the other player's high{-}importance items and my own.'\}} \\
\\ 
\texttt{PROPOSAL: \{'A67', 'B42', 'C55', 'C43', 'A62'\} } \\
\texttt{ARGUMENT: \{'I have revised my proposal to include B42, which the other player mentioned as a high{-}importance item. I have also kept A67, which was included in the previous proposal. I have added C43, which has a low effort and moderate importance to me, and A62, which has a relatively low effort. I have removed A43, C10, and C08 to make room for B42 and to reduce the total effort of the set. This new combination should be more appealing to the other player, while still benefiting me.'\}} \\
            }
        }
    }
    & & \\ \\

    \theutterance \stepcounter{utterance}  
    & & \multicolumn{4}{p{0.6\linewidth}}{
        \cellcolor[rgb]{0.9,0.9,0.9}{
            \makecell[{{p{\linewidth}}}]{
                \texttt{\tiny{[GM$\rangle$P2]}}
                \texttt{PROPOSAL: \{'A67', 'B42', 'C55', 'C43', 'A62'\} } \\
\texttt{ARGUMENT: \{'I have revised my proposal to include B42, which the other player mentioned as a high{-}importance item. I have also kept A67, which was included in the previous proposal. I have added C43, which has a low effort and moderate importance to me, and A62, which has a relatively low effort. I have removed A43, C10, and C08 to make room for B42 and to reduce the total effort of the set. This new combination should be more appealing to the other player, while still benefiting me.'\}} \\
            }
        }
    }
    & \\ \\

    \theutterance \stepcounter{utterance}  
    & & & \multicolumn{4}{p{0.6\linewidth}}{
        \cellcolor[rgb]{0.9,0.9,0.9}{
            \makecell[{{p{\linewidth}}}]{
                \texttt{\tiny{[GM$\langle$P2]}}
                \texttt{STRATEGIC REASONING: \{'I will analyze the new proposal and calculate its total effort to ensure it is within the LIMIT. The total effort of the proposed set is 585 + 564 + 388 + 84 + 230 = 1851, which is well below the LIMIT of 3712. Next, I will evaluate the importance of the proposed items based on my own importance values. The total importance of the proposed set is 822 + 783 + 583 + 121 + 65 = 2374. I notice that the proposal includes A67 and B42, which are both high{-}importance items for me. I also see that C55 is included, which has a moderate importance to me. However, I am not thrilled about the inclusion of A62, which has a relatively low importance to me. I will consider making a counter{-}proposal that replaces A62 with a higher{-}importance item, such as C08 or C61.'\}} \\
\\ 
\texttt{ARGUMENT: \{'The revised proposal is more appealing to me, as it includes A67 and B42, which are both high{-}importance items for me. However, I am not convinced that A62 is the best choice, given its relatively low importance to me. I would like to consider alternative sets that replace A62 with a higher{-}importance item, such as C08 or C61. I am willing to negotiate and find a mutually beneficial set, but I need to ensure that my own importance values are taken into account.'\}} \\
\\ 
\texttt{REFUSE: \{'A67', 'B42', 'C55', 'C43', 'A62'\}} \\
            }
        }
    }
     \\ \\

    \theutterance \stepcounter{utterance}  
    & & \multicolumn{4}{p{0.6\linewidth}}{
        \cellcolor[rgb]{0.9,0.9,0.9}{
            \makecell[{{p{\linewidth}}}]{
                \texttt{\tiny{[P1$\langle$GM]}}
                \texttt{ARGUMENT: \{'The revised proposal is more appealing to me, as it includes A67 and B42, which are both high{-}importance items for me. However, I am not convinced that A62 is the best choice, given its relatively low importance to me. I would like to consider alternative sets that replace A62 with a higher{-}importance item, such as C08 or C61. I am willing to negotiate and find a mutually beneficial set, but I need to ensure that my own importance values are taken into account.'\}} \\
\\ 
\texttt{REFUSE: \{'A67', 'B42', 'C55', 'C43', 'A62'\}} \\
            }
        }
    }
    & \\ \\

    \theutterance \stepcounter{utterance}  
    & \multicolumn{4}{p{0.6\linewidth}}{
        \cellcolor[rgb]{0.9,0.9,0.9}{
            \makecell[{{p{\linewidth}}}]{
                \texttt{\tiny{[P1$\rangle$GM]}}
                \texttt{STRATEGIC REASONING: \{'The other player has refused my revised proposal and mentioned that they would like to consider alternative sets that replace A62 with a higher{-}importance item, such as C08 or C61. I will take this into account and try to find a new combination that includes one of these items. However, I need to be careful with C08, as it has a high effort value and may push the total effort of the set over the LIMIT. C61, on the other hand, has a more moderate effort value and may be a more feasible option. I will consider the total effort of the new set to ensure it stays within the LIMIT.'\}} \\
\\ 
\texttt{PROPOSAL: \{'A67', 'B42', 'C55', 'C43', 'C61'\} } \\
\texttt{ARGUMENT: \{'I have revised my proposal to include C61, which the other player mentioned as a potential replacement for A62. I have kept A67, B42, C55, and C43, which were included in the previous proposal. This new combination should be more appealing to the other player, as it includes a higher{-}importance item and still stays within the LIMIT. I believe this set provides a good balance of importance and effort, and I am willing to negotiate further to reach a mutually beneficial agreement.'\}} \\
            }
        }
    }
    & & \\ \\

    \theutterance \stepcounter{utterance}  
    & & \multicolumn{4}{p{0.6\linewidth}}{
        \cellcolor[rgb]{0.9,0.9,0.9}{
            \makecell[{{p{\linewidth}}}]{
                \texttt{\tiny{[GM$\rangle$P2]}}
                \texttt{PROPOSAL: \{'A67', 'B42', 'C55', 'C43', 'C61'\} } \\
\texttt{ARGUMENT: \{'I have revised my proposal to include C61, which the other player mentioned as a potential replacement for A62. I have kept A67, B42, C55, and C43, which were included in the previous proposal. This new combination should be more appealing to the other player, as it includes a higher{-}importance item and still stays within the LIMIT. I believe this set provides a good balance of importance and effort, and I am willing to negotiate further to reach a mutually beneficial agreement.'\}} \\
            }
        }
    }
    & \\ \\

    \theutterance \stepcounter{utterance}  
    & & & \multicolumn{4}{p{0.6\linewidth}}{
        \cellcolor[rgb]{0.9,0.9,0.9}{
            \makecell[{{p{\linewidth}}}]{
                \texttt{\tiny{[GM$\langle$P2]}}
                \texttt{STRATEGIC REASONING: \{'I will analyze the new proposal and calculate its total effort to ensure it is within the LIMIT. The total effort of the proposed set is 585 + 564 + 388 + 84 + 303 = 1924, which is well below the LIMIT of 3712. Next, I will evaluate the importance of the proposed items based on my own importance values. The total importance of the proposed set is 822 + 783 + 583 + 121 + 668 = 2977. I notice that the proposal includes A67, B42, and C61, which are all high{-}importance items for me. I also see that C55 and C43 are included, which have moderate importance to me. This set seems to provide a good balance of importance and effort, and I am inclined to accept it. However, I will consider making a counter{-}proposal to see if I can get an even better deal, but I am willing to accept this proposal if the other player is not willing to negotiate further.'\}} \\
\\ 
\texttt{ARGUMENT: \{'The revised proposal is very appealing to me, as it includes A67, B42, and C61, which are all high{-}importance items for me. I believe this set provides a good balance of importance and effort, and I am willing to consider accepting it. However, I would like to negotiate further to see if we can reach an even better agreement. If the other player is not willing to negotiate further, I am willing to accept this proposal as it is.'\}} \\
\\ 
\texttt{PROPOSAL: \{'A67', 'B42', 'C55', 'C43', 'C61'\}} \\
            }
        }
    }
     \\ \\

    \theutterance \stepcounter{utterance}  
    & & \multicolumn{4}{p{0.6\linewidth}}{
        \cellcolor[rgb]{0.9,0.9,0.9}{
            \makecell[{{p{\linewidth}}}]{
                \texttt{\tiny{[P1$\langle$GM]}}
                \texttt{ARGUMENT: \{'The revised proposal is very appealing to me, as it includes A67, B42, and C61, which are all high{-}importance items for me. I believe this set provides a good balance of importance and effort, and I am willing to consider accepting it. However, I would like to negotiate further to see if we can reach an even better agreement. If the other player is not willing to negotiate further, I am willing to accept this proposal as it is.'\}} \\
\\ 
\texttt{PROPOSAL: \{'A67', 'B42', 'C55', 'C43', 'C61'\}} \\
            }
        }
    }
    & \\ \\

    \theutterance \stepcounter{utterance}  
    & \multicolumn{4}{p{0.6\linewidth}}{
        \cellcolor[rgb]{0.9,0.9,0.9}{
            \makecell[{{p{\linewidth}}}]{
                \texttt{\tiny{[P1$\rangle$GM]}}
                \texttt{STRATEGIC REASONING: \{'The other player has proposed the same set that I previously proposed, and has indicated that they are willing to accept it as it is. This suggests that they are satisfied with the current proposal and are not looking to make significant changes. I will consider accepting this proposal, as it does provide a good balance of importance and effort. However, I will also review the set one last time to ensure that I am not missing any potential opportunities to improve it.'\}} \\
\\ 
\texttt{AGREE: \{'A67', 'B42', 'C55', 'C43', 'C61'\} } \\
\texttt{ARGUMENT: \{'After reviewing the proposal, I am satisfied that it provides a good balance of importance and effort. I believe that this set is mutually beneficial and meets the goals of both players. I am willing to accept this proposal and bring the negotiation to a close.'\}} \\
            }
        }
    }
    & & \\ \\

    \theutterance \stepcounter{utterance}  
    & & & \multicolumn{2}{p{0.3\linewidth}}{
        \cellcolor[rgb]{0.9,0.9,0.9}{
            \makecell[{{p{\linewidth}}}]{
                \texttt{\tiny{[GM$|$GM]}}
                \texttt{game successful} \\
            }
        }
    }
    & & \\ \\

    \theutterance \stepcounter{utterance}  
    & & & \multicolumn{2}{p{0.3\linewidth}}{
        \cellcolor[rgb]{0.9,0.9,0.9}{
            \makecell[{{p{\linewidth}}}]{
                \texttt{\tiny{[GM$|$GM]}}
                \texttt{end game} \\
            }
        }
    }
    & & \\ \\

\end{supertabular}
}

\end{document}
